\documentclass[12pt,a4paper]{article}

% ── Packages ──────────────────────────────────────────────────────────────────
\usepackage[margin=1in]{geometry}
\usepackage{amsmath, amssymb}
\usepackage{listings}
\usepackage{xcolor}
\usepackage{parskip}
\usepackage{microtype}
\usepackage{hyperref}
\usepackage{booktabs}
\usepackage{caption}
\usepackage{setspace}
\usepackage{needspace}
\onehalfspacing

% ── Code listing style ─────────────────────────────────────────────────────────
\definecolor{codebg}{RGB}{245,245,245}
\definecolor{codeframe}{RGB}{200,200,200}
\definecolor{codecomment}{RGB}{100,100,100}
\definecolor{codekw}{RGB}{0,0,180}
\definecolor{codestr}{RGB}{163,21,21}

\lstset{
  language=Python,
  backgroundcolor=\color{codebg},
  frame=single,
  rulecolor=\color{codeframe},
  basicstyle=\ttfamily\small,
  keywordstyle=\color{codekw}\bfseries,
  stringstyle=\color{codestr},
  commentstyle=\color{codecomment}\itshape,
  showstringspaces=false,
  breaklines=true,
  numbers=left,
  numberstyle=\tiny\color{gray},
  tabsize=4,
  captionpos=b
}

% ── Title ──────────────────────────────────────────────────────────────────────
\title{
  \textbf{Face Recognition via Eigenfaces} \\[0.4em]
  \large A Mathematical and Code Walkthrough
}
\author{Aditya Maknikar \\ Machine Learning from Scratch}
\date{}

\begin{document}
\maketitle

\begin{center}
Mathematical Background: pg 1--4\\
Code Walkthrough: pg 5--11
\end{center}

% ==============================================================================
\needspace{5\baselineskip}
\section{Mathematical Background}
% ==============================================================================

\needspace{5\baselineskip}
\subsection{How Images Become Numbers}

Before any math can happen, the computer needs to turn a picture into something
it can actually work with. A grayscale image is stored as a rectangular grid of
pixels, where each pixel holds a single number representing brightness --- 0 is
black, 255 is white, and everything in between is a shade of grey.

Think of it like a spreadsheet: every row is a row of pixels, every column is a
column of pixels, and every cell contains a brightness value. A 112-by-92-pixel
image is literally just a $112 \times 92$ spreadsheet of numbers.

In the code, once an image is loaded it gets \emph{flattened} --- the 2D grid is
unrolled into a single long row of $L = 112 \times 92 = 10304$ numbers. One
image becomes one row in a larger matrix.

\needspace{5\baselineskip}
\subsection{The Dataset Matrix \texorpdfstring{$A$}{A}}

We have $M$ face images (after holding some out for testing, $M = 389$). After
flattening each image, we stack them on top of each other to form a matrix
$A$. This matrix has shape $M \times L$: each row is one image, each column
is one pixel position.

\[
A \in \mathbb{R}^{M \times L}, \quad M = 389,\; L = 10304
\]

Each row as one person's face, compressed into a single long list of
brightness values

\newpage
\needspace{5\baselineskip}
\subsection{Why We Subtract the Mean}

Before doing any analysis, we
subtract the mean face from every image in the dataset.

The mean vector $\mathbf{a}$ is computed by averaging across all $M$ images,
pixel by pixel. The result is a single vector of length $L$ --- the
``average face.'' We then subtract it from every row of $A$ to produce the
mean-centred matrix $C$:

\[
C = A - \mathbf{a}, \quad \text{where } \mathbf{a} = \frac{1}{M}\sum_{i=1}^{M} A_i
\]

We do this because we do not want our analysis to be dominated by the fact
that all faces share a common baseline of brightness and structure. If we skip
mean subtraction, the first principal component ends up describing ``faces in
general'' rather than ``the directions in which faces vary from each other.''
By subtracting the mean we force the analysis to focus on differences like the
variation in expression, lighting, pose, etc. that actually makes one
face distinct from another.

In the code this step is handled implicitly by \texttt{sklearn}'s
\texttt{PCA}, which stores the mean internally as \texttt{pca.mean\_} and
subtracts it for us. But conceptually this is the first thing that happens.

\needspace{5\baselineskip}
\subsection{The Covariance Matrix and Why It Is Enormous}

Once we have $C$, we want to understand the structure of the variation in the
data. The tool for that is the covariance matrix $S$:

\[
S = C^{\top} C
\]

Think of $S$ as a measure of how much any two pixel positions tend to vary
together across all the images. If pixel 500 and pixel 600 are both bright or
both dark in most images, their covariance is high. If one is bright when the
other is dark, the covariance is low (or negative). A problem we face is $S$ has shape $L \times L$, which means it is $10304 \times 10304$.
That is over 100 million entries. Computing the eigenvectors of a matrix that
large is not practical on ordinary hardware.

\newpage
\enlargethispage{2\baselineskip}
\needspace{5\baselineskip}
\subsection{The \texorpdfstring{$M \times M$}{M x M} Trick}

Instead of finding the eigenvectors of the $L \times L$ matrix $S = C^{\top}C$,
we can instead work with the much smaller $M \times M$ matrix:
\[
S' = C \cdot C^{\top}
\]
$S'$ has shape $M \times M = 389 \times 389$. Finding eigenvectors of
\emph{this} matrix is fast.

If $\mathbf{v}$ is an eigenvector of $C C^{\top}$ (the small matrix), then $C^{\top} \mathbf{v}$ is an eigenvector of $C^{\top} C$ (the big matrix), with the same eigenvalue. The two matrices share the same non-zero eigenvalues, and their eigenvectors are directly related. So we can compute eigenvectors of the small matrix and transform them into eigenvectors of the large one.

Stated formally: if $C C^{\top} \mathbf{v} = \lambda \mathbf{v}$, then:
\[
C^{\top} C \; (C^{\top} \mathbf{v}) = C^{\top}(C C^{\top} \mathbf{v}) = C^{\top}(\lambda \mathbf{v}) = \lambda (C^{\top} \mathbf{v})
\]
So $\mathbf{u} = C^{\top} \mathbf{v}$ is an eigenvector of $C^{\top}C$ with
eigenvalue $\lambda$. In practice, we then normalize $\mathbf{u}$ so it has
unit length. Scikit-learn handles this internally. when we call
\texttt{PCA().fit()}, this is the shortcut it uses . It does not
compute $C^{\top}C$ directly; it works with the smaller matrix to make the
computation tractable.

\needspace{5\baselineskip}
\subsection{Principal Components as Eigenfaces}

The eigenvectors of $S = C^{\top}C$ are the \emph{principal components} of the
face dataset. Because each eigenvector has length $L = 10304$, it can be
reshaped back into a $112 \times 92$ image. When you do that, the resulting
images look like blurry, ghostly faces. These are the eigenfaces.

The eigenvalues associated with each eigenvector tell us how much of the
variation in the dataset that eigenvector explains. The eigenvectors are ordered
from highest eigenvalue to lowest so the first eigenface captures the
single most important direction of variation, the second captures the next
most important, and so on.

We keep only the top $K$ eigenfaces (in the code, $K = 50$). The remaining ones
are discarded because they explain a very small fraction of the total variation.
Formally, we collect the top $K$ eigenvectors as rows of a matrix $F$ of shape
$K \times L$.

\newpage
\needspace{5\baselineskip}
\subsection{Representing Any Face as a Weight Vector}

Any face image can be
expressed as a weighted sum of the eigenfaces, plus the mean face. Eigen faces are similar to primary colours  and the weights say how much of each colour to use.

Formally, for a face image $\mathbf{x}$ (a row vector of length $L$), we
first subtract the mean to get the centred version
$\tilde{\mathbf{x}} = \mathbf{x} - \mathbf{a}$.
We then project this onto each eigenface by computing the dot product:
\[
    w_j = \mathbf{f}_j \cdot \tilde{\mathbf{x}}^{\top}, \quad j = 1, \ldots, K
\]
where $\mathbf{f}_j$ is the $j$-th eigenface (a row of $F$). The result is a
vector of $K$ weights:
\[
\mathbf{w} = F \cdot \tilde{\mathbf{x}}^{\top} \in \mathbb{R}^{K}
\]
This weight vector is a compact representation of the face. Two different
pictures of the same person should produce similar weight vectors, even if
the photos are not identical. A picture of a completely different person
should produce a noticeably different weight vector.

\needspace{5\baselineskip}
\subsection{Face Matching via Euclidean Distance}

To recognise a face, we compare its weight vector to the weight vectors of all
known faces and find the closest one. The distance metric used is the
Euclidean (L2) distance:

\[
d(\mathbf{w}_1, \mathbf{w}_2) = \|\mathbf{w}_1 - \mathbf{w}_2\|_2 =
\sqrt{\sum_{j=1}^{K}(w_{1j} - w_{2j})^2}
\]

The known face with the smallest distance to the query is the best match.
A threshold can optionally be applied: if the best match distance is above
the threshold, the face is declared ``unknown'' --- someone the model has
never seen.

\newpage

% ==============================================================================
\needspace{5\baselineskip}
\section{Code Walkthrough}
% ==============================================================================

The following walkthrough maps every block of code to the mathematical
concepts above. Each code block is shown alongside the explanation of what
it is actually doing.

\needspace{5\baselineskip}
\subsection{Loading the Images}

\vspace{0.5em}
\begin{minipage}{\linewidth}
\begin{lstlisting}[caption={Reading face images from a zip archive}]
import cv2
import zipfile
import numpy as np

faces = {}
with zipfile.ZipFile("attface.zip") as facezip:
    for filename in facezip.namelist():
        if not filename.endswith(".pgm"):
            continue  # not a face picture
        with facezip.open(filename) as image:
            faces[filename] = cv2.imdecode(
                np.frombuffer(image.read(), np.uint8),
                cv2.IMREAD_GRAYSCALE
            )
\end{lstlisting}
\end{minipage}

The dataset is the AT\&T Database of Faces: 40 people, 10 photos each, 400
images total. Each image is stored as a PGM file (grayscale, no colour
channels). The code iterates through the zip archive, reads each PGM file
into a byte buffer, and uses OpenCV's \texttt{cv2.imdecode} to turn that
buffer into a 2D NumPy array of pixel values.

After this block, \texttt{faces} is a Python dictionary mapping each filename
(e.g.\ \texttt{"s1/1.pgm"}) to its corresponding $112 \times 92$ NumPy array.
Nothing mathematical has happened yet --- this is purely data ingestion.

\needspace{5\baselineskip}
\subsection{Splitting into Training and Test Sets}

\begin{itemize}
  \item All 10 images of person 40 are held out. This person is entirely
        unknown to the model, so querying with their photo should produce a
        high distance (i.e.\ a failed recognition).
  \item One image of person 39 (\texttt{s39/10.pgm}) is held out. This person
        \emph{is} in the model, so querying with their photo should produce a
        low distance and a correct match.
\end{itemize}

\newpage



\vspace{0.5em}
\begin{minipage}{\linewidth}
\begin{lstlisting}[caption={Building the training matrix}]
facematrix = []
facelabel  = []

for key, val in faces.items():
    if key.startswith("s40/"):
        continue        # hold out entire person 40
    if key == "s39/10.pgm":
        continue        # hold out one image from person 39
    facematrix.append(val.flatten())
    facelabel.append(key.split("/")[0])

facematrix = np.array(facematrix)
\end{lstlisting}
\end{minipage}



For every image that is kept, \texttt{val.flatten()} converts the $112 \times 92$
2D array into a 1D vector of length $L = 10304$. This is the vectorisation
step described in Section~1.1. The final call \texttt{np.array(facematrix)}
stacks all those 1D vectors into the $M \times L$ matrix $A$ from Section~1.2.

The shape is therefore:
\[
A \in \mathbb{R}^{389 \times 10304}
\]
\texttt{facelabel} stores the class identifier for each row (e.g.\ \texttt{"s1"},
\texttt{"s2"}, \ldots, \texttt{"s39"}) so we know who each row corresponds to
when we look up a match later.

\needspace{5\baselineskip}
\subsection{Computing the Eigenfaces with PCA}

\vspace{0.5em}
\begin{minipage}{\linewidth}
\begin{lstlisting}[caption={Fitting PCA and extracting the top K eigenfaces}]
from sklearn.decomposition import PCA

pca = PCA().fit(facematrix)

n_components = 50
eigenfaces = pca.components_[:n_components]
\end{lstlisting}
\end{minipage}
\texttt{ PCA().fit(facematrix)} does the heavy lifting. Internally,sklearn performs the $M \times M$ trick described in Section~1.5 and returns
the principal components sorted by decreasing explained variance. We treat
this as a black box: we put our matrix $A$ in, and we get out the eigenvectors
of the covariance matrix $C^{\top}C$, ordered from most to least important.
\texttt{pca.components\_} is an array where each row is one principal component
vector one eigenface of length $L = 10304$. We keep only the first 50.
\newpage 
\[
F = \texttt{eigenfaces} \in \mathbb{R}^{50 \times 10304}
\]

\texttt{pca.mean\_} stores the mean face vector $\mathbf{a}$ of length $L$.
This is used in every subsequent computation to re-centre images before
projecting them.

The choice of $K = 50$ is a practical compromise. The explained variance
array printed from \texttt{pca.explained\_variance\_ratio\_} shows that the
first component alone explains about 17.8\% of total variance, the second
about 12.9\%, and the numbers drop off quickly after that. By 50 components,
we have captured the majority of the meaningful variation in the dataset
without carrying along hundreds of noisy, low-importance components.

\needspace{5\baselineskip}
\subsection{Visualising the Eigenfaces}

\vspace{0.5em}
\begin{minipage}{\linewidth}
\begin{lstlisting}[caption={Displaying the first 16 eigenfaces}]
import matplotlib.pyplot as plt

fig, axes = plt.subplots(4, 4, sharex=True, sharey=True, figsize=(8, 10))
for i in range(16):
    axes[i % 4][i // 4].imshow(
        eigenfaces[i].reshape(faceshape), cmap="gray"
    )
plt.show()
\end{lstlisting}
\end{minipage}

Each eigenface is stored as a flat vector of length 10304. To display it as
an image, \texttt{.reshape(faceshape)} converts it back into a $112 \times 92$
array. The eigenfaces look like blurry, ghostly face images because they are
literally the directions of maximum variation in face-pixel space. The first
few show broad lighting effects and general face structure. Further down the
list they pick up finer details like expression and hair.

\needspace{5\baselineskip}
\subsection{Computing Weight Vectors for All Training Faces}

\vspace{0.5em}
\begin{minipage}{\linewidth}
\begin{lstlisting}[caption={Projecting each training face onto the eigenface space}]
weights = eigenfaces @ (facematrix - pca.mean_).T
\end{lstlisting}
\end{minipage}

This single line of code implements the projection described in Section~1.6.
Breaking it down step by step:
\newpage
\begin{enumerate}
  \item \texttt{facematrix - pca.mean\_} subtracts the mean face $\mathbf{a}$
        from every row of $A$, producing the centred matrix $C$. This is an
        element-wise broadcast operation: NumPy automatically subtracts the
        $1 \times L$ mean vector from each of the $M$ rows.

  \item \texttt{.T} transposes $C$ from shape $M \times L$ to shape
        $L \times M$, so that each column is one mean-centred face.

  \item \texttt{eigenfaces @} performs matrix multiplication between $F$
        (shape $K \times L$) and $C^{\top}$ (shape $L \times M$). The
        \texttt{@} operator in Python is matrix multiplication (introduced in
        Python 3.5).
\end{enumerate}

The result is:

\[
W = F \cdot C^{\top} \in \mathbb{R}^{K \times M} = \mathbb{R}^{50 \times 389}
\]

Each column of $W$ is the weight vector for one training face. The $j$-th
entry in that column says how strongly the $j$-th eigenface contributes to
that particular image. Alternatively, each row of $W$ is the coefficient for
one eigenface across all 389 training images.

This is equivalent to the explicit loop below, which does the same
computation element-by-element but is much slower:

\vspace{0.5em}
\begin{minipage}{\linewidth}
\begin{lstlisting}[caption={Equivalent loop form (slower, for illustration only)}]
weights_loop = []
for i in range(facematrix.shape[0]):
    weight = []
    for j in range(n_components):
        w = eigenfaces[j] @ (facematrix[i] - pca.mean_)
        weight.append(w)
    weights_loop.append(weight)
\end{lstlisting}
\end{minipage}

The one-line matrix version and this loop produce identical results. The
matrix form is preferred in practice because it runs much faster.

\needspace{5\baselineskip}
\subsection{Recognising a Known Person}

\vspace{0.5em}
\begin{minipage}{\linewidth}
\begin{lstlisting}[caption={Query with a held-out image from a known class}]
query = faces["s39/10.pgm"].reshape(1, -1)
query_weight = eigenfaces @ (query - pca.mean_).T

euclidean_distance = np.linalg.norm(weights - query_weight, axis=0)
best_match        = np.argmin(euclidean_distance)

print("Best match %s with Euclidean distance %f"
      % (facelabel[best_match], euclidean_distance[best_match]))
\end{lstlisting}
\end{minipage}

The query image is person 39, photo 10 --- held out from the training set
but belonging to a person who \emph{is} in the model.

\texttt{query.reshape(1, -1)} converts the $112 \times 92$ image into a
$1 \times 10304$ row vector, matching the shape of the training rows in
\texttt{facematrix}.

\texttt{eigenfaces @ (query - pca.mean\_).T} repeats the projection from
Section~2.5, this time for just the one query image. The result
\texttt{query\_weight} has shape $K \times 1$ --- the weight vector for
the query.

\texttt{weights - query\_weight} broadcasts the subtraction across all 389
columns of \texttt{weights}. \texttt{np.linalg.norm(..., axis=0)} then
computes the Euclidean L2 distance between the query weight vector and each
of the 389 training weight vectors:

\[
d_i = \| \mathbf{w}_i - \mathbf{w}_{\text{query}} \|_2, \quad i = 1, \ldots, 389
\]

\texttt{np.argmin} returns the index of the smallest distance.
The result is:

\begin{center}
  \texttt{Best match s39 with Euclidean distance 1559.997137}
\end{center}

The model correctly identifies person 39. The relatively low distance
(compared to the unknown person test below) confirms that the query image
sits close to the other images of person 39 in the eigenface weight space.

\needspace{5\baselineskip}
\subsection{Attempting to Recognise an Unknown Person}

\vspace{0.5em}
\begin{minipage}{\linewidth}
\begin{lstlisting}[caption={Query with an image from a class never seen during training}]
query = faces["s40/1.pgm"].reshape(1, -1)
query_weight = eigenfaces @ (query - pca.mean_).T

euclidean_distance = np.linalg.norm(weights - query_weight, axis=0)
best_match        = np.argmin(euclidean_distance)

print("Best match %s with Euclidean distance %f"
      % (facelabel[best_match], euclidean_distance[best_match]))
\end{lstlisting}
\end{minipage}

This block is identical to the previous one, but the query image is person 40 ---
someone the model has never seen. The model cannot return the correct answer
because person 40's weight vector simply is not in the training set. The
result is:

\begin{center}
  \texttt{Best match s5 with Euclidean distance 2690.209330}
\end{center}

Two things to note. First, the best-match distance is noticeably higher
(2690 vs.\ 1560). This is the signal that tells us the query is likely an
unknown person. Second, the model still returns \emph{some} best match (s5),
because \texttt{np.argmin} always returns the closest thing it can find,
even if the closest is still quite far.

This motivates the threshold idea: in a real system, you would choose a
cutoff value $\tau$ and say: if $d_{\min} < \tau$, report the best match;
if $d_{\min} \geq \tau$, report ``unknown.'' Choosing $\tau$ well is an
empirical problem that requires a validation set.

\needspace{5\baselineskip}
\subsection{Generating a New Face from Random Weights}

\vspace{0.5em}
\begin{minipage}{\linewidth}
\begin{lstlisting}[caption={Constructing a face from randomly sampled weights}]
random_weights = np.random.randn(n_components) * weights.std()
newface = random_weights @ eigenfaces + pca.mean_

fig, axes = plt.subplots(1, 2, sharex=True, sharey=True, figsize=(8, 6))
axes[0].imshow(pca.mean_.reshape(faceshape), cmap="gray")
axes[0].set_title("Mean face")
axes[1].imshow(newface.reshape(faceshape), cmap="gray")
axes[1].set_title("Random face")
plt.show()
\end{lstlisting}
\end{minipage}

This block demonstrates the reconstruction direction: instead of going from
image to weights, we go from weights to image.

\texttt{np.random.randn(n\_components)} samples $K = 50$ values from a
standard normal distribution. Multiplying by \texttt{weights.std()} scales
them to roughly the same range as real face weights, so the result does not
look completely unrecognisable.

\texttt{random\_weights @ eigenfaces} computes a weighted sum of the eigenfaces:
\[
\tilde{\mathbf{z}} = \sum_{j=1}^{K} w_j \mathbf{f}_j \in \mathbb{R}^{L}
\]
Adding \texttt{pca.mean\_} restores the mean face, giving the final
reconstructed image:
\[
\mathbf{z} = \tilde{\mathbf{z}} + \mathbf{a}
\]
The result is a face-like image that does not correspond to any real person.
It looks plausible because the eigenfaces encode real statistical structure
from the training data --- the result is constrained to the ``face-shaped''
region of pixel space that the PCA model has learned.

\needspace{5\baselineskip}
\subsection{Accuracy and Limitations}

The eigenface method is surprisingly effective given how simple the underlying
math is. Testing under controlled conditions has found recognition rates around
96\% when lighting is consistent, dropping to about 85\% with changes in
head orientation, and down to 64\% when subject distance varies. The core
problem is that PCA is a linear method operating on raw pixel values. Changes
in lighting, scale, or pose produce large movements in pixel space that do
not correspond to changes in identity. A zoom-in or zoom-out, for instance,
shifts pixel values in a way the model has no way to factor out.


\end{document}
